\documentclass[lang=en]{sjtureport}
\usepackage{amsmath}
\usepackage{booktabs}
\usepackage{siunitx}

\title{A Measurement Report on Queueing Delay}
\author{Jordan Blake}
\subject{Systems Engineering Laboratory}
\keywords{queueing, latency, reproducibility}

\begin{document}
\maketitle

\chapter{Introduction}
This report studies how burst size changes the waiting time in a single-server
queue. The purpose is not to claim a universal performance law, but to show a
complete path from a question to a reproducible measurement and a bounded
conclusion.

\chapter{Experimental design}
Requests arrive in batches at a fixed interval. The service time is sampled
from a documented distribution, and the queue is reset before each run. For a
stable system with arrival rate $\lambda$ and service rate $\mu$, the traffic
intensity is
\begin{equation}
  \rho=\frac{\lambda}{\mu},\qquad 0\leq\rho<1.
\end{equation}
Runs with $\rho\geq1$ are reported as saturated rather than treated as normal
steady-state observations.

\chapter{Results}
\begin{table}[ht]
\centering
\caption{Illustrative queue measurements.}
\begin{tabular}{@{}lrrr@{}}
\toprule
Batch size & Requests & Median wait (ms) & 95th percentile (ms)\\
\midrule
1 & 600 & 4.2 & 8.7\\
4 & 600 & 9.5 & 18.3\\
8 & 600 & 19.1 & 41.6\\
\bottomrule
\end{tabular}
\end{table}

The tail grows faster than the median as the workload becomes burstier. Before
comparing implementations, repeat each row with the same seed, warm-up rule,
hardware, and clock source.

\chapter{Discussion}
The measurements support a workload-specific association between burst size
and queueing delay. They do not isolate scheduler overhead, garbage
collection, or network variance. Record those factors explicitly in a future
run and include confidence intervals around both central and tail estimates.

\chapter{Conclusion}
For the tested service rate, smaller batches keep the tail shorter. The result
is useful as a diagnostic baseline; capacity planning should use production
traces and a saturation test.

\appendix
\chapter{Run checklist}
\begin{enumerate}
  \item Record the software revision and machine configuration.
  \item Declare warm-up, sampling, and outlier handling before looking at results.
  \item Archive raw observations with the generated report.
\end{enumerate}
\end{document}
